\begin{statementbox}
	\textbf{\textcolor{CStatement}{条件}}
	\begin{description}[style=nextline, leftmargin=2.8em, labelsep=0.5em, itemsep=0.35em]
		\item[\textbf{\textcolor{CStatement}{条件 1（函数二次）：}}]
			二次函数 $f(x)=ax^2+bx+c$，$a\neq0$，对称轴 $x_0=-\frac{b}{2a}$ 固定。
		\item[\textbf{\textcolor{CStatement}{条件 2（区间动）：}}]
			区间 $[m,n]$ 端点可变，满足 $m<n$。
	\end{description}
	\textbf{\textcolor{CStatement}{结论}}
	\begin{description}[style=nextline, leftmargin=2.8em, labelsep=0.5em, itemsep=0.45em]
		\item[\textbf{\textcolor{CStatement}{结论一（开口向上）：}}]
			\textbf{\textcolor{CStatement}{直观描述：}}
			当 $a>0$ 时，最小值取决于对称轴是否在区间内，最大值始终是两端点较大者。
			\[
				f_{\min}=
				\begin{cases}
					f(x_0), & \text{若 } x_0\in[m,n],\\ \min\{f(m),f(n)\}, & \text{若 } x_0\notin[m,n],
			\end{cases}
				\qquad f_{\max}= \max\{f(m),f(n)\}.
			\]
		\item[\textbf{\textcolor{CStatement}{结论二（开口向下）：}}]
			\textbf{\textcolor{CStatement}{直观描述：}}
			当 $a<0$ 时，最大值取决于对称轴是否在区间内，最小值始终是两端点较小者。
			\[
				f_{\max}=
				\begin{cases}
					f(x_0), & \text{若 } x_0\in[m,n],\\ \max\{f(m),f(n)\}, & \text{若 } x_0\notin[m,n],
			\end{cases}
				\qquad f_{\min}= \min\{f(m),f(n)\}.
			\]
		\item[\textbf{\textcolor{CStatement}{常见变形（对称性应用）：}}]
			\textbf{\textcolor{CStatement}{直观描述：}}
			利用二次函数对称性，若区间关于 $x_0$ 对称，则 $f(x_0)$ 必为最值（开口向上时最小，开口向下时最大），且两端点函数值相等。
	\end{description}

\end{statementbox}
